\chapter{Empalme material entre la estructura única y sus cinco resoluciones}
\label{chap:pdf2-generacion-simultanea-hmt}

\begingroup
\footnotesize
% CORTE_2_NO_DISGREGACION — bloque de empalme PDF1 -> PDF2
% Congelado después del PDF1 compilado el 2026-08-19.
% Owners:
% 04b  8ba02f246c18a3bd2ff83ffdb08874c816a7f59c688b6fab6af6eb2f9647071f
% 04b1 25d6dee3875cbdd3252f793212448c17c4b5b2d3c5d787e995598d18fa76d5cf
% 04b2 d5de9b5379d4f95ea3d35e63c9ccc2ab741b5c750c01409af6f13526073b874f
% 04b3 ac9ca321b610e2c0560a32ec6ae1555814c67c587e152096d0c48d6abaa0bd4b

\section*{Empalme matemático material PDF1--PDF2}

La entrada única es el estado APP--TRIT--TPK completo. Antes de toda
evaluación arquimediana se recorren
\[
 0\longrightarrow10,\qquad 0\longrightarrow1000,
 \qquad 0\longrightarrow\infty,
\]
con los nueve huecos internos \(1,\ldots,9\); dos posiciones producen APP,
TRIT orienta sus hojas y TPK prolonga extensión, fase, carry, memoria,
frontera, ruta y supervivencia. Ni el continuo ni \(\mathbb R\) son entradas.

Para una emisión superviviente
\[
 \mathfrak e_\alpha=
 (\alpha,\operatorname{Ext}_{\gamma_\alpha},\operatorname{or}(\alpha),
 \kappa(\alpha),\partial\alpha,\gamma_\alpha,
 m(\gamma_\alpha;m_0),\operatorname{Led}(\gamma_\alpha))
\]
las cinco características son acciones simultáneas sobre ese mismo
generador:
\begin{align*}
 \mathcal O_{\rm sp}(\mathfrak e_\alpha)
  &=(U_\alpha,\sigma_x,\sigma_y,A_\alpha,\eta_\alpha),\\
 A_\alpha&=P_yU_\alpha P_x+Q_yU_\alpha Q_x,
 &\eta_\alpha&=Q_yU_\alpha P_x+P_yU_\alpha Q_x,\\
 \mathcal O_{\rm sol}(\mathfrak e_\alpha)
  &=(b_\alpha,b_\alpha^+,b_\alpha^-;
  M^\pm\mapsto M^\pm+b_\alpha^\pm,
  D\mapsto D+b_\alpha,H\mapsto H+|b_\alpha|),\\
 \mathcal O_{\rm gau}(\mathfrak e_\alpha)
  &=\{(\operatorname{Cell}_{\rm TPK}(\alpha),q,B,W_c(q)):
      q\in Q_{\rm inc}(t\alpha)\},\\
 \mathcal O_{\rm coh}(\mathfrak e_\alpha)
  &=(I_y,J_y,\kappa_{y,N},\delta_y,C_{y,N},
      F_\alpha:C_{y,N}\to C_{x,N}),\\
 \mathcal O_{\rm det}(\mathfrak e_\alpha)
  &=(\Psi_\alpha,K_\alpha,\mathbf1_\alpha,z_\pm(\alpha),h_*(\alpha),
  \operatorname{Det}C_{y,N},H_*C_{y,N},\operatorname{Smith}C_{y,N},
  \operatorname{Boc}C_{y,N}).
\end{align*}
Estas expresiones son vistas dependientes de una sola acción
\(\mathcal O(\mathfrak e_\alpha)\); no definen cinco dominios, factores o
ramas. Comparten literalmente fuente, destino, ruta, ledger, orientación,
carry, memoria y frontera. Sus identidades de acoplamiento son
\[
 A_\alpha^*A_\alpha+\eta_\alpha^*\eta_\alpha=U_\alpha^*U_\alpha,
 \qquad E b_f^\pm=b_c^\pm+\kappa^{\rm can},
\]
\[
 Q_{\rm inc}=C^2(\operatorname{Cell}_{\rm TPK};\mathbf F_3),
 \qquad sB=0,\qquad W_c(q+Ba)=W_c(q),
\]
\[
 \partial h_*=z_--z_+,qquad
 K_{\beta\alpha}=K_\beta+\Psi_\beta K_\alpha.
\]

La monodromía es la conexión nonádica real:
\[
 \Gamma_{9,k}=\operatorname{Hol}_{C_9}(\nabla_{\rm TPK}),
 \qquad g(k+9)=g(k),
 \qquad q_{\rm mem}(k+9)=q_{\rm mem}(k)+1,
\]
sin retorno de estado. Para una historia
\(h=(\varepsilon_{k+1},\ldots,\varepsilon_N)\), las trece coordenadas se
generan por el fold
\[
 r_{k,N}(h)=
 \operatorname{Upd}_{\varepsilon_N}\circ\cdots\circ
 \operatorname{Upd}_{\varepsilon_{k+1}}(r^0_{k}),
\]
y el terminal devuelve toda la multisección, sin selector:
\[
 \mathfrak T_{k,N}(\widehat x_k)
 =\{r_{k,N}(h):h\in\mathscr H_{k,N}(\widehat x_k)\}.
\]
La naturalidad es una única igualdad sobre prefijos generados,
\[
 u_{\ell,k}^{\mathcal O}\mathcal O_{\le\ell}(\gamma)
 =\mathcal O_{\le k}(\gamma)
 =\operatorname{map}(\mathcal O)
   \bigl(u_{\ell,k}^{\mathfrak e}E_{\le\ell}(\gamma)\bigr),
\]
que pasa al sistema inverso de terminales. Sólo entonces se obtiene la
estructura discreta completa del continuo; la evaluación en \(\mathbb R\) es
una publicación posterior.

\subsection*{Mapa de consumo en PDF2}

\begin{center}
\small
\begin{tabularx}{\textwidth}{@{}p{.18\textwidth}X p{.30\textwidth}@{}}
\toprule
Símbolo consumido & Contenido estructural íntegro de PDF1 &
Entrelazador que construye PDF2 \\
\midrule
\(\mathfrak G_{\rm sp}\) &
\(U,\sigma,P,Q,A,\eta\), hojas, carry, terminal, naturalidad y lector
firmado interno. &
\(\Sigma_R^{\rm sp},J_\pm,K_R^{\rm sp},\Xi,P,B_W\) y la identificación
con la forma de Weil. \\
\(\mathfrak G_{\rm sol}\) &
\(b,b^\pm,\kappa^{\rm can}\), memoria, micro/shell ortogonales, columna
terminal, Gram, Stein y telescopía sin doble conteo. &
Mapa source-first \(J_T^{\rm sol}\), \(\Lambda_{M,T}\), \(I_{M,J}\),
factorización uniforme y cota sólo en datos. \\
\(\mathfrak G_{\rm gau}\) &
\(\operatorname{Cell}_{\rm TPK},Q_{\rm inc},B,W_c\), cuatro direcciones,
seis celdas, ocho caras, expectativas y terminal común. &
Interfaz gauge física, semigrupo requerido, sector físico, Yang--Mills y
brecha de masa. \\
\(\mathfrak G_{\rm coh}\) &
Puertos, \(\kappa,\delta,C\), bar/coend, carry, supervivencia y cono de
holonomía. &
\(R_{\rm Hdg}\), \(\mathrm{ch}_{\rm loc}\), clase algebraica y publicación
Hodge. \\
\(\mathfrak G_{\rm det}\) &
\(\Psi,K,z_\pm,h_*\), AW, línea determinante, homología, Smith y
Bockstein, conservados sin exigir aciclicidad como entrada. &
Kato--FERS, PT, comparador local--\newline global, Tamagawa y publicación BSD. \\
\bottomrule
\end{tabularx}
\end{center}

Los objetos de la tercera columna son target-specific: no deben retrotraerse
como entradas de PDF1. PDF2 los construye sobre el objeto estructural completo
de la segunda columna y conserva sus trece coordenadas, terminal y genealogía.

\endgroup

\section{Tipado unitario del objeto consumido}

El bloque congelado anterior no se añade a un constructor gráfico de cinco
componentes. Lo sustituye. Para cada profundidad \(k\), sea
\(\mathcal E_k\) el estado APP--TRIT--TPK enriquecido con todas sus emisiones
supervivientes, y sea \(\mathcal Z_k\) el codominio de la única acción
correlativa. La flecha rectora es
\begin{equation}\label{eq:pdf2-constructor-unico-operaciones}
 \mathcal O_k:\mathcal E_k\longrightarrow\mathcal Z_k,
 \qquad
 \widehat x_k\longmapsto
 \mathcal O\bigl(\{\mathfrak e_\alpha:
 \alpha\in\operatorname{Ch}_k(\widehat x_k)\}\bigr).
\end{equation}
Su valor conserva simultáneamente las coordenadas
\(\mathrm{sp},\mathrm{sol},\mathrm{gau},\mathrm{coh},\mathrm{det}\)
del mismo estado y del mismo ledger. No es una tupla de factores
independientes ni el grafo tautológico de cinco mapas previamente separados.

Los truncamientos del estado y de la acción se escriben
\[
 u^{\mathfrak e}_{\ell,k}:\mathcal E_\ell\to\mathcal E_k,
 \qquad
 u^{\mathcal O}_{\ell,k}:\mathcal Z_\ell\to\mathcal Z_k,
 \qquad \ell\ge k,
\]
y satisfacen una sola naturalidad,
\begin{equation}\label{eq:pdf2-naturalidad-accion-unica}
 u^{\mathcal O}_{\ell,k}\circ\mathcal O_\ell
 =\mathcal O_k\circ u^{\mathfrak e}_{\ell,k}.
\end{equation}
Por ello la acción pasa al límite sin seleccionar una historia:
\[
 \mathcal O_\infty:
 \mathcal C_{\rm cont}^{\rm disc}:=\varprojlim_k\mathcal E_k
 \longrightarrow
 \mathcal Z_\infty:=\varprojlim_k\mathcal Z_k.
\]
Las notaciones usadas en los cinco capítulos son vistas intrínsecas de ese
único valor:
\begin{equation}\label{eq:pdf2-operacion-infinita-tipificada}
 \mathfrak G_T^{\rm int}
 :=\operatorname{View}_T\!\left(
 \mathcal O_\infty(\mathcal C_{\rm cont}^{\rm disc})\right),
 \qquad
 T\in\{\mathrm{sp},\mathrm{sol},\mathrm{gau},\mathrm{coh},\mathrm{det}\}.
\end{equation}
Aquí \(\operatorname{View}_T\) sólo expone las coordenadas requeridas por una
aplicación; no crea una rama, un factor, un subdominio ni un selector. La
estructura de partida de cada resolución es siempre el valor completo de
\(\mathcal O_\infty\), con las otras cuatro características presentes.

\begin{falsifier}[Orden causal y no disgregación]
Falla el empalme si se usa \(\mathbb R\) o el continuo como entrada de
\(\mathcal O_k\), si se reemplaza \(\mathcal O_k\) por cinco constructores
independientes, si se selecciona una sola historia terminal o si una vista
\(\operatorname{View}_T\) se confunde con el objeto completo. En cualquiera
de esos casos la conclusión del capítulo de aplicación deja de transferirse.
\end{falsifier}
